% ============================================================
%  chapters/04_data.tex
% ============================================================

\section{Data}
\label{sec:data}

\subsection{Nigeria Education Data Survey (NEDS)}

The primary data for this paper come from the 2010 Nigeria Education Data
Survey (NEDS). The survey was conducted by the National Population Commission
(NPC), the Federal Ministry of Education (FMOE), and the Universal Basic
Education Commission (UBEC), with technical assistance from RTI
International and joint funding from the United States Agency for
International Development (USAID) and the UK Department for International
Development (DFID). The objectives of the 2010 NEDS include
\citep{neds2010}:

\begin{itemize}
    \item Providing data on the schooling status of Nigerian children of
          basic education age, including factors influencing whether children
          ever enroll in school and why students drop out
    \item Measuring parent attitudes toward schooling quality and providing
          an understanding of the attitudes that shape willingness to send
          children to school
    \item Measuring the frequency of student absenteeism and reasons for
          missing school
    \item Providing data that allow for trend analysis and state-level
          comparisons
\end{itemize}

The survey targeted children aged 4--16 with a focus on the household-level
factors affecting education consumption decisions. Data collection was
conducted across all 36 states and the Federal Capital Territory (Abuja)
from April to July 2010. In total, 26,934 households were interviewed,
yielding information on 71,567 children.

\subsection{Uppsala Conflict Data Program (UCDP)}

The violence data were collected by Uppsala University as part of the
Uppsala Conflict Database Program. The program's purpose is to provide
comprehensive event-level data on all organized violence since 1989. An
event is defined as ``any incidence where there is a use of armed force by
an organized actor against another organized actor or civilians, which
results in at least one direct death at a specific location and for a
specific temporal duration'' \citep{sundberg2010}.\footnote{See
Appendix~\ref{app:violence_defs} for complete UCDP definitions of key
terms including organized actor, direct death, and event types.} Uppsala
University maintains accuracy in the dataset through a four-step data
collection process involving systematic searches of major newswires and
databases, review of secondary materials, consultation with regional
experts, and final review by program directors.\footnote{See
Appendix~\ref{app:ucdp_process} for a description of the four-step data
collection process.}

The dataset records date(s), event type, violence type, location details
(town, province, and latitude/longitude coordinates), actors involved, and
deaths for each side as well as civilian and unknown deaths. Events are
classified as single-day, summary, or continuous events. Violence types are
categorized as state-based, non-state, or one-sided violence. For this
study, I use all events recorded in Nigeria from 1999 to 2010.

\subsection{Linking the Datasets}

In order to link the household survey data with the violence data, I matched
records using state and LGA identifiers present in both datasets. Some
observations in the violence dataset did not contain LGA-level information
and were therefore excluded from the matched sample. Summary statistics for
the key variables used in the analysis are presented in
Table~\ref{tab:summary_stats}.

% ---- Placeholder for summary statistics table ---------------
% Once you have your data cleaned and merged, export a summary
% statistics table from Python/Stata and include it here.
% See the LaTeX tables guide in main.tex for formatting.

\begin{table}[H]
    \centering
    \caption{Summary Statistics}
    \label{tab:summary_stats}
    \begin{threeparttable}
        \begin{tabular}{lcccc}
            \toprule
            Variable & Mean & Std. Dev. & Min & Max \\
            \midrule
            \multicolumn{5}{l}{\textit{Panel A: Outcomes}} \\[0.3em]
            Currently enrolled (\textit{attend}) & & & 0 & 1 \\
            Years of schooling (\textit{highest\_enrolled}) & & & 0 & \\[0.5em]
            \multicolumn{5}{l}{\textit{Panel B: Violence Exposure}} \\[0.3em]
            Violence: pre-birth ($-$17 to $-$9 months) & & & 0 & 1 \\
            Violence: in utero (0–9 months pre-birth) & & & 0 & 1 \\
            Violence: early childhood (ages 0–3) & & & 0 & 1 \\
            Violence: preschool (ages 3–6) & & & 0 & 1 \\
            Violence: primary school (ages 6–12) & & & 0 & 1 \\
            Violence: junior secondary (ages 12–15) & & & 0 & 1 \\
            Violence: senior secondary (ages 15–18) & & & 0 & 1 \\[0.5em]
            \multicolumn{5}{l}{\textit{Panel C: Individual Controls}} \\[0.3em]
            Male & & & 0 & 1 \\
            Age & & & & \\
            Urban & & & 0 & 1 \\
            Wealth quintile & & & 1 & 5 \\
            Distance to primary school (km) & & & & \\
            Mother attended school & & & 0 & 1 \\
            Father attended school & & & 0 & 1 \\
            \midrule
            Observations & \multicolumn{4}{c}{[N]} \\
            \bottomrule
        \end{tabular}
        \begin{tablenotes}
            \small
            \item \textit{Notes:} Summary statistics for the analysis
            sample. Violence indicators are binary variables equal to 1 if
            the child's LGA experienced at least one UCDP-recorded violent
            event during the specified age window.
            Violence exposure is assigned based on the child's LGA of
            residence and the start date of the first UCDP event in that
            LGA. See Section~\ref{sec:methodology} for details on variable
            construction.
        \end{tablenotes}
    \end{threeparttable}
\end{table}
